在上一讲中，我们讨论了静电平衡导体的性质：导体内部电场为零，导体本身是等势体，表面是等势面，
多余电荷只能分布在表面。对于一个已经达到静电平衡的导体，电势和净电荷是描述其状态的两个重要物理量。
若附近还有别的导体，一个导体所带电荷改变时，不仅自身电势会改变，还会通过静电感应改变其他导体表面的
电荷分布和电势。因此，导体组的静电问题本质上是在研究各导体的电荷、电势与几何结构之间的关系。

本讲先用电容描述这种关系，再讨论电容器的连接与储能，随后把结论推广到任意导体组，最后研究电介质在
电场中的极化。以下定义都默认导体的形状、相对位置及周围介质已经固定；只有在这些条件不变时，电荷与
电势之间的线性关系才可以由一个固定的电容或一组电容系数表示。

\subsection*{孤立导体的电容}

设空间中有若干导体，除其中一个以外，其余导体都接地，电势保持为零。使唯一未接地的导体带上净电荷
$Q$，静电平衡后它的电势为 $\varphi$。由静电场的线性叠加和唯一性定理可知：如果把 $Q$ 同时放大
$k$ 倍，空间中全部电场和电势也放大 $k$ 倍，因而 $Q/\varphi$ 不随电荷量改变。定义该导体在当前
几何环境中的单极电容为
\begin{equation}
  \boxed{C=\frac{Q}{\varphi}},\qquad Q=C\varphi.
\end{equation}
这里的“孤立”不是说宇宙中只有这一个导体，而是说只把它作为非零电势的导体，其余导体接地。电容取决于
导体的大小、形状、位置及周围介质，却与它在给定线性范围内带了多少电荷无关。国际单位制中电容的单位是
法拉（F），名称来自 Faraday；实际导体的电容通常远小于一法拉，常用 $\si{\micro F}$、$\si{nF}$、
$\si{pF}$。

半径为 $R$ 的孤立导体球是最容易计算的例子。球外电势与位于球心的点电荷相同，取无穷远电势为零，
\begin{equation}
  \varphi(R)=\frac{Q}{4\pi\varepsilon_0R},
  \qquad
  \boxed{C_{\rm sphere}=4\pi\varepsilon_0R}.
\end{equation}
结果与球所带电量无关，并且正比于球半径。球越大，在升高到同一电势时能够容纳的电荷越多。

给一个原来不带电的孤立导体逐渐充电。当它已经带有电荷 $q$ 时，再搬入微小电荷 $\dd q$ 所做的功是
$\varphi(q)\dd q=(q/C)\dd q$。从零充到 $Q$，系统的静电能为
\begin{equation}
  W=\int_0^Q\frac{q}{C}\dd q
  =\boxed{\frac{Q^2}{2C}=\frac12Q\varphi=\frac12C\varphi^2}.
\end{equation}
这些能量并不是“藏在电荷里”，而是以电场能的形式分布在空间中。真空静电场的能量密度为
$u_E=\varepsilon_0E^2/2\geq0$，所以静电系统的场能不会为负。

\subsection*{双极电容器与典型几何}

实际储存电荷时更常用两块彼此绝缘的导体。令两导体分别带 $+Q$ 与 $-Q$，电势为
$\varphi_+$、$\varphi_-$，两极电势差
$U=\varphi_+-\varphi_-$。只要几何与介质不变，$Q$ 与 $U$ 成正比，双极电容器的电容定义为
\begin{equation}
  \boxed{C=\frac{Q}{U}},\qquad Q=CU.
\end{equation}
孤立导体的 $Q/\varphi$ 与双极电容器的 $Q/U$ 形式相同，但物理边界条件不同，使用时不能混淆。

\arbitraryCapacitorDiagram

电容器不要求极板平行，也不要求形状规则。它的电容由两极的大小、形状、相对位置及介质共同决定。
当导体形状具有足够高的对称性时，可以先用 Gauss 定律求电场，再积分求电势差，最后取 $C=Q/U$。

对于面积为 $S$、间距为 $d$ 的平行板电容器，若极板线度远大于 $d$，中部电场近似均匀：
\begin{equation}
  E=\frac{Q}{\varepsilon_0S},\qquad
  U=Ed=\frac{Qd}{\varepsilon_0S},
  \qquad
  \boxed{C=\frac{\varepsilon_0S}{d}}.
\end{equation}
真实极板边缘的电场会弯曲，称为边缘效应。极板越大、间距越小，均匀场区域在全部场区中所占比例越大，
理想公式越准确。边缘场使实际电容通常略大于忽略边缘效应得到的数值。

同心球形电容器的内、外半径为 $R_1<R_2$。在两球面之间
\begin{equation}
  E(r)=\frac{Q}{4\pi\varepsilon_0r^2},
  \qquad
  U=\int_{R_1}^{R_2}E(r)\dd r,
\end{equation}
因此
\begin{equation}
  \boxed{C_{\rm sph}=4\pi\varepsilon_0\frac{R_1R_2}{R_2-R_1}}.
\end{equation}
令 $R_2\to\infty$，便退化为半径 $R_1$ 的孤立导体球。

同轴圆柱电容器内、外半径为 $a<b$，有效长度为 $l$，并且 $l$ 足够大，可忽略两端效应。若内柱带
$+Q$，两极之间距轴线 $r$ 处的场强为
\begin{equation}
  E(r)=\frac{Q}{2\pi\varepsilon_0lr}.
\end{equation}
由
\begin{equation}
  U=\int_a^bE(r)\dd r
   =\frac{Q}{2\pi\varepsilon_0l}\ln\frac ba
\end{equation}
得到
\begin{equation}
  \boxed{C_{\rm cyl}=\frac{2\pi\varepsilon_0l}{\ln(b/a)}}.
\end{equation}
圆柱电容正比于长度。外半径固定而内半径减小时，靠近细内导体处的场会迅速增强，这也说明“无限细带电线”
只是理论极限，真实导线半径会影响电势差和电容。

\capacitorGeometryDiagram

\subsection*{电容器的连接}

多个电容器接入电路后，可以用一个等效电容描述从外部端口看到的储荷效果。并联电容器的两端分别连接到
同一对节点，故各电容器的电势差都等于 $U$。第 $i$ 个电容器带电 $Q_i=C_iU$，电源搬运的总电荷为
\begin{equation}
  Q=\sum_iQ_i=\left(\sum_iC_i\right)U,
\end{equation}
所以
\begin{equation}
  \boxed{C_{\parallel}=\sum_iC_i}.
\end{equation}

串联电容器的中间连接导体与外界绝缘。若它初始不带净电荷，充电过程中电荷只能在这个导体内部重新分布，
净电荷始终为零，因而各串联电容器所带电荷量的绝对值相同。总电势差为
\begin{equation}
  U=\sum_iU_i=Q\sum_i\frac1{C_i},
\end{equation}
于是
\begin{equation}
  \boxed{\frac1{C_{\rm series}}=\sum_i\frac1{C_i}}.
\end{equation}
并联使等效电容增大，串联使等效电容减小。电容的串并联规律与电阻恰好相反，却与理想弹簧的等效关系在
形式上相似；真正的判断依据是各元件“电压相同”还是“电荷相同”。

\seriesParallelDiagram

\subsection*{电容储能与能量密度}

容量为 $C$ 的电容器最初不带电。当两极已经带有 $\pm q$ 时，电势差为 $u=q/C$，再把 $\dd q$ 从
负极搬到正极需要外力做功
\begin{equation}
  \dd W=u\dd q=\frac qC\dd q.
\end{equation}
从零充到 $Q$ 后，电容器储存的电场能为
\begin{equation}
  \boxed{W=\frac{Q^2}{2C}=\frac12QU=\frac12CU^2}.
\end{equation}
这个推导不依赖极板形状。对于平行板电容器，把上式除以板间体积 $Sd$，可得真空中的能量密度
\begin{equation}
  \boxed{u_E=\frac12\varepsilon_0E^2}.
\end{equation}
在线性、各向同性介质中，能量密度推广为 $u_E=\vect E\cdot\vect D/2=\varepsilon E^2/2$。

\subsection*{任意导体组与 Green 倒易定理}

孤立导体和双极电容器都只是导体组的特殊情况。设同一组固定导体有两个静电平衡状态：第一状态中第 $i$
个导体的电荷与电势为 $Q_i,\varphi_i$，第二状态为 $Q_i',\varphi_i'$。由 Green 第二恒等式可得
\begin{equation}
  \boxed{\sum_iQ_i\varphi_i'=\sum_iQ_i'\varphi_i}.
\end{equation}
它说明：把第一状态各导体的电荷乘以第二状态相应导体的电势后求和，与交换两个状态所得结果相同。
具有这种“交换激励与响应而结果不变”结构的结论称为倒易定理。它与静电场方程的线性、边界条件的唯一性
以及场能的二次型结构密切相关。

\reciprocitySphereDiagram

例如，取一个导体球和球外距球心为 $r$ 的点 $P$。状态 I 中球带单位电荷，$P$ 点无点电荷；球面电荷
均匀分布，它在 $P$ 点产生电势 $1/(4\pi\varepsilon_0r)$。状态 II 中球原来不带电，在 $P$ 点放置
电荷 $q$，静电感应使球面电荷重新分布，球仍为等势体，设其电势为 $\varphi$。对这两个状态应用倒易定理，
可立即得到
\begin{equation}
  \boxed{\varphi=\frac{q}{4\pi\varepsilon_0r}}.
\end{equation}
虽然状态 II 的球面电荷分布并不均匀，但球的电势无需逐点积分就能确定。

\reciprocityConductorsDiagram

再考虑两个任意形状导体。两种状态分别记为
$(Q_1,Q_2;\varphi_1,\varphi_2)$ 与
$(Q_1',Q_2';\varphi_1',\varphi_2')$，则
\begin{equation}
  Q_1\varphi_1'+Q_2\varphi_2'
  =Q_1'\varphi_1+Q_2'\varphi_2.
\end{equation}
由此可以导出多种有用结论。例如，在两种实验中轮流令一个导体带相同电荷、另一个保持中性，则一个导体
在另一个导体处产生的电势具有交换对称性；轮流令一个导体接地、另一个处于相同电势，接地导体上的总感应
电荷也具有交换对称性。把电荷写成各导体电势的线性函数
\begin{equation}
  Q_i=\sum_jC_{ij}\varphi_j,
\end{equation}
倒易定理给出
\begin{equation}
  \boxed{C_{ij}=C_{ji}}.
\end{equation}
因此，多导体静电感应并不是每一对导体各有一套互不相关的经验参数，而可以用一个对称电容系数矩阵统一描述。

\subsection*{电介质与相对介电常数}

真空电容器的两极之间没有物质。实际电容器往往在电极间填入橡胶、玻璃、纸、石英等绝缘材料，这类物质
称为电介质。实验表明：用某种均匀电介质完全填满电容器后，电容总会增大；更重要的是，对同一种介质，
无论原电容器形状如何，在适用的线性范围内，电容都近似按同一倍数增大。这个倍数称为相对介电常数
$\varepsilon_r$：
\begin{equation}
  \boxed{C=\varepsilon_rC_0},
\end{equation}
其中 $C_0$ 是相同几何结构在真空中的电容。常见被动介质的 $\varepsilon_r>1$。

\dielectricExperimentDiagram

\begin{center}
\small
\begin{tabular}{lcc lcc}
\toprule
材料 & 相对介电常数 $\varepsilon_r$ & 击穿场强/$(\si{V.m^{-1}})$ &
材料 & 相对介电常数 $\varepsilon_r$ & 击穿场强/$(\si{V.m^{-1}})$\\
\midrule
干燥空气 & $1.00059$ & $3\times10^6$ & 水 & $80$ & --\\
石英 & $3.78$ & $24\times10^6$ & 橡胶 & $6.7$ & $12\times10^6$\\
纸 & $3.7$ & $16\times10^6$ & 尼龙 & $3.4$ & $14\times10^6$\\
玻璃 & $5.6$ & $14\times10^6$ & 聚苯乙烯 & $2.56$ & $24\times10^6$\\
\bottomrule
\end{tabular}

{\small\color{TextGray}表 32.1\quad 室温下典型介质的相对介电常数和击穿电场（典型值）}
\end{center}

表中的击穿场强表示介质仍能保持绝缘的典型上限。超过这一数量级后，束缚电荷可能被大量释放，介质发生
电击穿，不能再用线性绝缘体模型处理。材料的介电常数与击穿场强还会受到温度、频率、纯度和制备状态影响，
表中数值只适合作数量级估计。

若保持电容器自由电荷 $Q$ 不变，完全插入介质后 $C$ 增大为 $\varepsilon_rC_0$，电势差
$U=Q/C$ 降低到原来的 $1/\varepsilon_r$，场强也相应减小。这说明介质内部出现了一个通常与外场反向的
附加电场。它来源于介质中束缚正、负电荷在外场下的微小相对位移。

\subsection*{介质极化与电位移}

经典图像中，绝缘体由大量整体中性的原子或分子组成，电子不能像金属自由电子那样在整个物体内迁移。
外电场作用时，原子中正、负电荷中心发生微小分离，形成诱导电偶极子；本来具有永久偶极矩的分子也会在
外场下出现取向偏好。两种机制都使大量微观偶极矩在统计平均后产生宏观极化。

单位体积内的总电偶极矩定义为极化强度
\begin{equation}
  \boxed{\vect P=\frac{\dd\vect p}{\dd V}}.
\end{equation}
若每个分子的平均偶极矩为 $\vect p$、数密度为 $N$，均匀极化时 $\vect P=N\vect p$。介质内部相邻偶极子
的正负电荷大体抵消，但在边界上会留下束缚面电荷；非均匀极化还会产生束缚体电荷：
\begin{equation}
  \boxed{\sigma_b=\vect P\cdot\uvec n},
  \qquad
  \boxed{\rho_b=-\nabla\cdot\vect P}.
\end{equation}

\polarizationDiagram

自由电荷与束缚电荷都产生真实电场。为了在介质问题中把由电极等带来的自由电荷单独列出，定义电位移矢量
\begin{equation}
  \boxed{\vect D=\varepsilon_0\vect E+\vect P}.
\end{equation}
于是 Gauss 定律可写为
\begin{equation}
  \boxed{\oint_S\vect D\cdot\dd\vect S=Q_{\rm free}},
  \qquad
  \nabla\cdot\vect D=\rho_{\rm free}.
\end{equation}
右侧只计自由电荷，极化形成的束缚电荷已经包含在 $\vect P$ 中。$\vect D$ 是重新组织电荷贡献后引入的
辅助场量，并不是另一种独立于电场的力。

“电位移”这个名称保留了电磁学发展早期的历史痕迹。十九世纪的机械以太模型曾把电场类比为某种连续介质
受力后的形变，把真空响应与普通物质的极化相加。以太图像后来被放弃，但由线性场方程整理出的
$\vect D=\varepsilon_0\vect E+\vect P$ 仍然是处理介质问题的有效定义。学习这一历史背景的意义在于区分：
公式可以保留，而最初解释公式的机械模型未必正确。

当场强不太大、介质各向同性且响应近似即时，极化强度可近似与电场成正比：
\begin{equation}
  \vect P=\varepsilon_0\chi_e\vect E.
\end{equation}
因此
\begin{equation}
  \boxed{\vect D=\varepsilon\vect E
  =\varepsilon_r\varepsilon_0\vect E},
  \qquad \varepsilon_r=1+\chi_e.
\end{equation}
从微观角度，若单个原子的诱导偶极矩满足 $\vect p=\alpha\vect E_{\rm loc}$，$\alpha$ 称为原子极化率。
宏观介电常数与 $N\alpha$ 有关；在最简单的各向同性模型中还要考虑邻近偶极子对局域场的修正，因此不能
总把 $\vect E_{\rm loc}$ 直接等同于宏观平均场 $\vect E$。测量介电常数曾是由宏观实验反推原子结构的重要途径。

\subsection*{介质边界、部分填充与受力}

在两种介质的交界面上作极薄柱形 Gauss 面和极窄矩形回路，可得静电边界条件
\begin{equation}
  \boxed{D_{2n}-D_{1n}=\sigma_f},
  \qquad
  \boxed{E_{2t}=E_{1t}}.
\end{equation}
只有自由面电荷使 $\vect D$ 的法向分量发生跳变；静电场的切向分量连续。界面没有自由电荷时
$D_{1n}=D_{2n}$，但若两侧介电常数不同，$E_n$ 一般不连续，电场线通过界面时会改变方向。

介质沿电场方向分层时，各层自由电荷面密度相同，$D_n$ 相同，而总电压等于各层电压之和，等效为电容串联：
\begin{equation}
  \frac1C=\frac{d_1}{\varepsilon_1S}+\frac{d_2}{\varepsilon_2S}+\cdots.
\end{equation}
介质若横向分割极板面积，各部分电势差相同、总电荷相加，等效为并联：
\begin{equation}
  C=\frac{\varepsilon_1S_1}{d}+\frac{\varepsilon_2S_2}{d}+\cdots.
\end{equation}
判断串联或并联不能只看图形，而要看各部分是电荷相同还是电压相同。

介质板插入电容器通常受到向内部的力，因为插入会使电容增大。电容器孤立、自由电荷固定时
\begin{equation}
  W_Q=\frac{Q^2}{2C},
  \qquad
  F_x=-\left(\frac{\partial W_Q}{\partial x}\right)_Q
  =\frac{Q^2}{2C^2}\frac{\partial C}{\partial x}.
\end{equation}
若电容器接恒压电源，必须把电源做功一起计入。电容器与电源合成系统的有效势能变化为
$\dd(W_C-W_s)=-U^2\dd C/2$，所以
\begin{equation}
  \boxed{F_x=\frac12U^2\frac{\partial C}{\partial x}}.
\end{equation}
两种边界条件给出的力方向一致，但推导中保持不变的量不同。

\subsection*{充电过程中的能量去向}

把未充电电容 $C$ 通过电阻 $R$ 接到理想电压源 $U_0$，最终电源搬运电荷 $Q=CU_0$，做功
$W_s=U_0Q=CU_0^2$。电容最终储能只有 $CU_0^2/2$，另一半在电阻中转化为热：
\begin{equation}
  Q_R=\int_0^\infty I^2R\dd t=\frac12CU_0^2.
\end{equation}
这一半能量与 $R$ 的具体数值无关，$R$ 只改变充电快慢。理想化地令 $R\to0$ 时，集总模型出现无限大
初始电流；真实电路中的寄生电感、电磁辐射和有限传播速度会参与能量转移，不能据此断言能量损耗消失。
串联电阻充电时
\begin{equation}
  U_C(t)=U_0\left(1-\e^{-t/(RC)}\right),
  \qquad
  I(t)=\frac{U_0}{R}\e^{-t/(RC)}.
\end{equation}
$RC$ 是时间常数。经过一个时间常数，电容电压达到终值的 $1-\e^{-1}$，电流降为初值的 $\e^{-1}$。

\subsection*{从边值问题理解电容}

电容的定义看上去只是一个比值，但它实际上概括了完整的静电边值问题。以孤立导体为例，给定导体表面的
电势以后，导体外部的电势必须满足 Laplace 方程，并在导体表面取给定常数、在无穷远趋于零。唯一性定理
保证满足这些条件的解只有一个。若把表面电势由 $\varphi$ 改成 $k\varphi$，原解整体乘以 $k$ 就满足新的
边界条件；于是电场、表面电荷密度和总电荷也全部乘以 $k$。这正是 $Q=C\varphi$ 成立的根本原因。

在这个构造中，其他接地导体虽然电势为零，却不等于“没有作用”。它们表面会出现感应电荷并改变空间电场，
因而也会改变未接地导体的电容。若未接地导体带正电，邻近接地导体朝向它的一侧通常出现负感应电荷。
接地意味着大地可以供给或吸收电子，所以这些导体的净电荷不必保持为零。把接地导体移近，通常会使给定
电荷造成的电势降低，从而使 $C=Q/\varphi$ 增大。这说明“一个导体的电容”仍然包含周围边界的信息。

双极电容器也由同样的线性原理决定。两极带等量异号电荷时，正极发出的电场线最终终止在负极，理想情况下
电场主要局限在两极之间。电容器储存的不是“正电荷”或“负电荷”本身，而是把正负电荷分离后形成的电场能。
如果保持两极形状不变，只把它们靠近，给定 $Q$ 造成的电势差会减小，因此电容增大；如果增大正对面积，
给定总电荷分摊到更大表面上，场强与电势差降低，电容同样增大。平行板公式中的 $S$ 和 $d$ 依赖关系由此
获得直观解释，而不只是代数结果。

平行板近似还包含一个容易被忽略的尺度条件。若极板边长与间距同量级，边缘弯曲的电场占据很大比例，不能
把全部场区都视为匀强场；此时 $C=\varepsilon_0S/d$ 只能给出粗略数量级。将极板边长同时放大而保持间距
不变，边缘区域的宽度仍约为 $d$，但中部匀强区域所占比例越来越大，近似才逐渐变好。竞赛题若强调
$l\gg d$，就是在要求忽略这部分端部修正。

圆柱电容的积分也展示了几何如何进入电容。圆柱 Gauss 面的面积正比于 $r$，所以场强按 $1/r$ 衰减；
电势差是 $1/r$ 从 $a$ 到 $b$ 的积分，因而出现 $\ln(b/a)$。如果只记结论而忽略这一过程，容易把它与球形
电容的 $1/R_1-1/R_2$ 混淆。三种典型电容器可用统一步骤处理：确定对称性，选取 Gauss 面求 $E$，沿场线
积分求 $U$，最后用 $C=Q/U$ 消去 $Q$。

\subsection*{连接、充电和储能的完整推理}

并联电容的关键是端点相同。达到静电平衡后，理想导线上的电势处处相同，因此每只电容的两端电压都等于
外部端口电压。电源向整个并联系统搬运的电荷是各电容正极板电荷之和，故等效电容相加。串联电容的关键
则是中间节点孤立：这个节点由相邻两块极板和连接导线构成，外部没有导电通路。它可以在内部重新分配正负
电荷，却不能改变净电荷；初始中性时，相邻两电容上的电荷量绝对值必然相同。把每段电压 $Q/C_i$ 相加，
便得到倒数相加的规律。

这种推理比直接背公式更可靠。例如串联电容的中间节点若预先带有净电荷，各电容所带电荷量便不再简单相等；
若电容器之间不是理想导线而连接了别的元件，节点条件也要重新分析。所谓“串联”或“并联”并不是由图形
看起来像一条线还是两条支路决定，而是由电势约束和电荷守恒决定。

电容储能中的因子 $1/2$ 来源于充电过程中电压从零逐渐增加。若错误地把最终电压 $U$ 乘以全部电荷 $Q$，
会得到 $QU$，比真实储能大一倍。正确做法是把每一小份电荷在当时电压 $u=q/C$ 下搬运所需的功积分起来。
也可以在 $q-u$ 图上理解：充电过程对应从原点到 $(Q,U)$ 的直线，储能等于直线下的三角形面积
$QU/2$。对于线性介质，$Q$ 与 $U$ 成正比；若介质具有强非线性或滞后，储能必须沿真实充电路径积分，
不能直接套用 $CU^2/2$。

把电场能密度在空间积分，应当得到与充电做功相同的总能量。以平行板电容为例，板间场近似均匀，
\[
 \int u_E\dd V=\frac12\varepsilon_0E^2Sd
 =\frac12\frac{\varepsilon_0S}{d}(Ed)^2
 =\frac12CU^2.
\]
这一等价性说明，集中参数 $C$ 的储能公式与分布在空间中的场能描述的是同一个物理事实。遇到介质只填充
部分空间、场分布不均匀的情形，直接积分场能有时比先求等效电容更方便。

电源给电容充电时还要区分“电源做功”和“电容最终储能”。理想恒压源每搬运 $\dd q$ 都做功
$U_0\dd q$，所以总功为 $U_0Q$；而电容电压从零升到 $U_0$，平均电压只有 $U_0/2$，最终储能是
$U_0Q/2$。差额必须进入电阻发热、辐射或电磁场的暂态过程。把电阻取得很小只会让过程更快，不会凭空改变
初态与终态之间的能量差。

\subsection*{倒易定理的三种典型用法}

Green 倒易定理的力量在于，它允许把一个难求的响应换成另一个容易计算的响应。应用时必须保证两个状态中
导体的形状、位置和介质完全相同，仅改变各导体的电荷或电势。若几何也改变，两个状态不再属于同一个线性
边值问题，不能直接使用倒易关系。

第一种典型用法是交换“点电荷”和“测量电势的位置”。在导体球外放点电荷时，球面感应电荷分布很复杂；
但把球改成带单位电荷、去掉外部点电荷后，球面电荷均匀，任意外点电势都容易写出。倒易定理把第二种状态
的已知电势转换成第一种状态中整个球的电势，从而避开求感应电荷密度的繁重计算。

第二种用法是交换两个导体。状态 I 中让导体 1 带电 $Q$、导体 2 中性；状态 II 中让导体 2 带相同电荷
$Q$、导体 1 中性。代入倒易式可知，第一状态中导体 2 的电势等于第二状态中导体 1 的电势。这个结论与
两导体的形状是否相同、距离多远都无关，只要求几何在交换状态时没有改变。

第三种用法涉及接地感应电荷。状态 I 中把导体 1 提升到单位电势而让导体 2 接地，状态 II 中交换两导体。
倒易定理说明，两次实验中接地导体上出现的总感应电荷相同。用矩阵语言表述，就是电容系数的互易性
$C_{12}=C_{21}$。这一性质使复杂导体系统所需的独立系数数量大幅减少，也是后来网络互易定理的静电学原型。

倒易定理还可以和能量观点联系。导体组总静电能可写成
\[
 W=\frac12\sum_iQ_i\varphi_i
   =\frac12\sum_{ij}C_{ij}\varphi_i\varphi_j.
\]
交换两个电势变量不应改变同一个能量函数的混合二阶变化，因此系数矩阵必须对称。这给出了倒易关系的另一
种直观来源：互易性并非偶然的代数巧合，而是线性系统存在单值能量函数的结果。

\subsection*{从实验现象到介质极化}

插入电介质的实验可以有两种做法。若电容器先充电后与电源断开，自由电荷 $Q$ 不变，插入介质后电压表读数
降低；由 $C=Q/U$ 可知电容增大。若电容器始终与恒压电源相连，电压不变，插入介质时电源继续向极板输送
电荷；由 $Q=CU$ 可知储荷量增加。两种实验观测的是同一事实，只是外部保持不变的量不同。

同一种介质填满不同形状的电容器时，电容近似按同一比例增加，这一点非常关键。它说明介质作用不能只用
某种特定极板几何来解释，而应由材料自身的响应参数描述。相对介电常数正是这样的宏观参数。表 32.1 中
不同物质的数值相差很大，水的相对介电常数远大于干燥空气，表明分子结构和极化机制会显著影响宏观响应。
击穿场强则描述另一件事：介质在多强的场下还能维持绝缘。介电常数大并不意味着击穿场强必然大，设计真实
电容器时必须同时考虑储能能力和绝缘可靠性。

在保持自由电荷不变的实验中，插入介质后宏观场减小。微观上，外场使原子或分子的正负电荷中心分离，
形成与外场同向的电偶极矩；大量偶极子在介质两端留下束缚电荷。束缚电荷产生的场在介质内部通常与原外场
反向，于是合场减弱。这里“偶极矩与外场同向”和“偶极子附加场与外场反向”并不矛盾：前者描述单个偶极
矩的方向，后者描述极化介质端面束缚电荷在材料内部产生的退极化场。

均匀极化时介质内部的束缚体电荷为零，但表面束缚电荷一般不为零；非均匀极化时 $\rho_b=-\nabla\cdot\vect P$
也会贡献电场。因此，不能把极化简单理解为“只有表面有电荷”。具体是表面项还是体积项占主导，取决于
$\vect P$ 的空间分布和边界形状。

引入 $\vect D$ 的主要好处，是在已知自由电荷而束缚电荷难以逐个计算时直接使用 Gauss 定律。例如平行板
间完全填充均匀线性介质，极板自由电荷面密度为 $\sigma_f$，则 $D=\sigma_f$，进一步由
$D=\varepsilon E$ 得 $E=\sigma_f/\varepsilon$。相对于真空，场强减小 $\varepsilon_r$ 倍，电势差也减小
相同倍数，电容因而增大 $\varepsilon_r$ 倍。这个推导把实验规律、极化图像和宏观场方程连接起来。

不过 $\vect D$ 不能替代 $\vect E$ 描述带电粒子的受力。试探电荷受到的仍是 $q\vect E$，不是
$q\vect D$。在介质界面上，$\vect D$ 的法向跳变由自由面电荷决定，$\vect E$ 的切向分量由静电场保守性
保证连续；只有把两个条件结合起来，才能正确判断电场线如何折转。

\subsection*{线性模型的适用范围}

$\vect P=\varepsilon_0\chi_e\vect E$ 是近似的材料关系。场强过大时，极化可能不再与场成正比；极性分子
的取向还会受到热运动影响，因此介电常数可能随温度变化。交变电场频率很高时，较慢的极化机制来不及响应，
介电常数会随频率变化并出现能量损耗。各向异性晶体中，$\vect P$ 也未必与 $\vect E$ 平行，此时介电常数
应当用张量而不是一个标量表示。

原子极化率 $\alpha$ 是微观量，介电常数是宏观量。若简单地认为每个原子所见局域场就是宏观平均场，便有
$\vect P=N\alpha\vect E$；但邻近偶极子也会改变原子所在位置的场，所以精确关系需要局域场修正。经典的
Clausius--Mossotti 关系写成
\begin{equation}
  \boxed{\frac{\varepsilon_r-1}{\varepsilon_r+2}
  =\frac{N\alpha}{3\varepsilon_0}},
\end{equation}
它展示了如何由宏观介电常数估计微观极化率。该式适用于由近似独立、各向同性分子组成的稀疏介质；对于
强关联固体、极性液体或接近相变的材料，不能机械套用。


\subsection*{重点知识}

\begin{enumerate}[label=\textbf{32.\arabic*},leftmargin=4.2em]
  \item 双极电容器的电容为两极电荷量的绝对值与电势差之比，$C=Q/U$。
  \item 孤立导体的电容为导体电荷量与相对无穷远电势之比，可视为它与无穷远之间的电容。
  \item 电容器储能 $W=Q^2/(2C)=QU/2=CU^2/2$，能量以电场形式存在。
  \item 绝缘体称为电介质。电介质会极化，但一般不会像导体那样形成静电屏蔽和内部等势状态。
  \item 电位移满足 $\vect D=\varepsilon_0\vect E+\vect P$；其 Gauss 定律只以自由电荷为源。
\end{enumerate}


% 以下为早期压缩稿，保留在源码中仅供修订对照，不进入成品。
\iffalse
导体达到静电平衡后，每块连通导体都具有确定的电势与净电荷。固定所有导体的形状、位置以及
周围介质后，静电场方程是线性的，因此导体的电荷与电势之间也呈线性关系。对单个孤立导体或
只关心两极电势差的双极系统，这种线性关系就由一个标量——电容——来概括。

\subsection*{孤立导体与双极电容器}

若一个导体携带电荷 $Q$，相对于无穷远的电势为 $\varphi$，其余导体均接地，则定义单极电容
\begin{equation}
  \boxed{C=\frac{Q}{\varphi}}.
\end{equation}
电容由几何形状、相对位置和介质决定，与 $Q$ 的大小无关。半径为 $R$ 的孤立导体球满足
$\varphi=Q/(4\pi\varepsilon_0R)$，故
\begin{equation}
  \boxed{C_{\rm sphere}=4\pi\varepsilon_0R}.
\end{equation}

两块导体分别带等量异号电荷 $+Q$、$-Q$，两极电势差为 $U$ 时，双极电容器的电容定义为
\begin{equation}
  \boxed{C=\frac QU}.
\end{equation}
这一定义不要求极板平行；只要结构固定且介质在线性范围内，$Q$ 与 $U$ 就成正比。

平行板电容器在极板面积 $S$ 远大于边缘尺度 $d^2$ 时，可忽略边缘场。板间
$E=Q/(\varepsilon_0S)$、$U=Ed$，所以
\begin{equation}
  \boxed{C=\frac{\varepsilon_0S}{d}}.
\end{equation}
同心球形电容器内、外半径为 $R_1<R_2$，同轴圆柱电容器内、外半径为 $a<b$、有效长度为 $l$，
则
\begin{equation}
  \boxed{C_{\rm sph}=4\pi\varepsilon_0\frac{R_1R_2}{R_2-R_1}},
  \qquad
  \boxed{C_{\rm cyl}=\frac{2\pi\varepsilon_0l}{\ln(b/a)}}.
\end{equation}

\capacitorGeometryDiagram

\subsection*{连接与储能}

并联电容器两端电势差相同，总电荷是各支路电荷之和：
\begin{equation}
  \boxed{C_{\parallel}=\sum_iC_i}.
\end{equation}
串联电容器的中间连接导体若初始中性，其净电荷始终为零，因而各电容器所带电荷量的绝对值
相同；总电压是各段电压之和：
\begin{equation}
  \boxed{\frac1{C_{\rm series}}=\sum_i\frac1{C_i}}.
\end{equation}

\seriesParallelDiagram

给电容器增加微小电荷 $\dd q$ 时，需要克服当时的电势差 $u=q/C$，外力输入
$\dd W=u\dd q$。从零电荷充到 $Q$，电场储能为
\begin{equation}
  \boxed{W=\int_0^Q\frac qC\dd q
  =\frac{Q^2}{2C}=\frac12QU=\frac12CU^2}.
\end{equation}
真空中的能量密度为 $u_E=\varepsilon_0E^2/2$；在线性、各向同性介质中推广为
\begin{equation}
  \boxed{u_E=\frac12\vect E\cdot\vect D=\frac12\varepsilon E^2}.
\end{equation}

\subsection*{导体组的倒易关系}

同一组固定导体可处于两种静电平衡状态。第一态各导体的电荷、电势为 $Q_i,\varphi_i$，第二态
为 $Q_i',\varphi_i'$。由 Green 第二恒等式，可得静电倒易定理
\begin{equation}
  \boxed{\sum_iQ_i\varphi_i'=\sum_iQ_i'\varphi_i}.
\end{equation}
把电荷写成电势的线性函数
\begin{equation}
  Q_i=\sum_jC_{ij}\varphi_j,
\end{equation}
倒易定理给出 $C_{ij}=C_{ji}$。这意味着“导体 $j$ 升高单位电势在导体 $i$ 上引起的电荷”与
交换 $i,j$ 后完全相同。多导体静电感应并不需要为每一对导体分别建立互不相关的参数，电容系数
矩阵具有对称性。

\subsection*{电介质、极化与电位移}

绝缘体置于外电场后，束缚的正负电荷发生微小相对位移，原子或分子形成电偶极子。单位体积内的
总电偶极矩定义为极化强度
\begin{equation}
  \boxed{\vect P=\frac{\dd\vect p}{\dd V}}.
\end{equation}
均匀介质中，体内相邻偶极子的电荷相互抵消，但介质表面会留下束缚面电荷
\begin{equation}
  \boxed{\sigma_b=\vect P\cdot\uvec n}.
\end{equation}
非均匀极化还会产生束缚体电荷 $\rho_b=-\nabla\cdot\vect P$。

\polarizationDiagram

定义电位移矢量
\begin{equation}
  \boxed{\vect D=\varepsilon_0\vect E+\vect P}.
\end{equation}
把束缚电荷并入 $\vect P$ 后，Gauss 定律写成
\begin{equation}
  \boxed{\oint_S\vect D\cdot\dd\vect S=Q_{\rm free}},
  \qquad \nabla\cdot\vect D=\rho_{\rm free}.
\end{equation}
对线性、各向同性、均匀介质，
\begin{equation}
  \vect P=\varepsilon_0\chi_e\vect E,
  \qquad
  \boxed{\vect D=\varepsilon\vect E
  =\varepsilon_r\varepsilon_0\vect E},
  \qquad \varepsilon_r=1+\chi_e.
\end{equation}
同一电容器完全填充介质后，若几何不变，电容增大为
\begin{equation}
  \boxed{C=\varepsilon_rC_0}.
\end{equation}
电介质不具备导体的静电屏蔽性质：极化只能削弱或改变电场，不能让介质内部在一般情况下变成
等势体。

\fi

\Needspace{0.42\textheight}
\begin{origexamplebox}{例题 32.1【统统算出来】}
求下列三种电容器的电容：
\begin{enumerate}[label=(\arabic*),leftmargin=2.2em]
  \item 面积 $S=l^2$、间距 $d$ 的平行板电容器，$l\gg d$；
  \item 内、外半径为 $R_1,R_2$ 的同心球形电容器，$R_1<R_2$；
  \item 内、外半径为 $a,b$、长度为 $l$ 的同轴圆柱形电容器，$l\gg b-a$。
\end{enumerate}
\end{origexamplebox}

\begin{xesolution}
三题都让内、外电极分别带 $+Q$、$-Q$，先由 Gauss 定律求两极间电场，再积分求电压。

\textbf{(1) 平行板。}忽略边缘效应，板间场为
\[
  E=\frac{Q}{\varepsilon_0S},\qquad
  U=Ed=\frac{Qd}{\varepsilon_0S}.
\]
因此
\begin{equation}
  \boxed{C_{\rm pp}=\frac QU=\frac{\varepsilon_0S}{d}}.
\end{equation}

\textbf{(2) 同心球。}在 $R_1<r<R_2$ 内，
\[
  E(r)4\pi r^2=\frac Q{\varepsilon_0},
  \qquad E(r)=\frac{Q}{4\pi\varepsilon_0r^2}.
\]
故
\[
  U=\int_{R_1}^{R_2}E(r)\dd r
  =\frac{Q}{4\pi\varepsilon_0}\left(\frac1{R_1}-\frac1{R_2}\right),
\]
从而
\begin{equation}
  \boxed{C_{\rm sph}=4\pi\varepsilon_0
  \frac{R_1R_2}{R_2-R_1}}.
\end{equation}
令 $R_2\to\infty$，得到孤立球的 $C=4\pi\varepsilon_0R_1$。

\textbf{(3) 同轴圆柱。}柱面 Gauss 面给出
\[
  E(r)2\pi rl=\frac Q{\varepsilon_0},
  \qquad E(r)=\frac{Q}{2\pi\varepsilon_0lr}.
\]
所以
\[
  U=\int_a^bE(r)\dd r=\frac{Q}{2\pi\varepsilon_0l}\ln\frac ba,
\]
以及
\begin{equation}
  \boxed{C_{\rm cyl}=\frac{2\pi\varepsilon_0l}{\ln(b/a)}}.
\end{equation}
三种推导中的近似只用于忽略开口或极板边缘的非均匀场；在理想同心球中不存在边缘近似。
\end{xesolution}

\Needspace{0.38\textheight}
\begin{origexamplebox}{例题 32.2【两个导体球】}
两个半径均为 $R$ 的导体球紧挨着并相互导通。求它们作为一个导体时的单极电容。
\end{origexamplebox}

\begin{xesolution}
两球导通后具有相同电势 $V$。这一几何可用双球坐标求解，也可用逐次反演产生的无穷像电荷
序列。对两个等半径球，设球心距为 $L>2R$，定义
\[
  \cosh\alpha=\frac{L}{2R}.
\]
两个球处于同一电势 $V$ 时，自电容系数与互电容系数的级数交错组合，总电荷可写成
\begin{equation}
  Q_{\rm tot}=8\pi\varepsilon_0RV\sinh\alpha
  \sum_{n=1}^{\infty}\frac{(-1)^{n+1}}{\sinh(n\alpha)}.
\end{equation}
因此
\[
  C(L)=8\pi\varepsilon_0R\sinh\alpha
  \sum_{n=1}^{\infty}\frac{(-1)^{n+1}}{\sinh(n\alpha)}.
\]
相切极限 $L\to2R^+$ 对应 $\alpha\to0^+$。此时
\[
  \frac{\sinh\alpha}{\sinh(n\alpha)}\longrightarrow\frac1n,
\]
故级数趋于交错调和级数 $1-1/2+1/3-\cdots=\ln2$，得到
\begin{equation}
  \boxed{C=8\pi\varepsilon_0R\ln2}.
\end{equation}
它约为 $1.386$ 个孤立单球电容 $4\pi\varepsilon_0R$，小于相距无穷远时两个单球电容之和
$8\pi\varepsilon_0R$。原因是两球接近且同电势时，接触区附近的同号电荷被排斥到外侧，总电势在
相同总电荷下升高。

\touchingSpheresDiagram

图 32.12 中沿球心连线画出的点列表示双球坐标或像电荷级数中的等效源位置。越靠近接触点，
相邻像电荷的位置越密，说明接触区的边界条件需要越来越高阶的项才能准确满足；这也是相切极限
不能简单把两个孤立球电容直接相加的原因。图中的点不是导体内部真实存在的点电荷，而是计算外部
电势时的数学等效源。

结果的两个极限提供了重要校验。若两球相距很远，每个球几乎不受另一个球影响，总电容趋近
$2(4\pi\varepsilon_0R)=8\pi\varepsilon_0R$；相切时得到
$8\pi\varepsilon_0R\ln2$，因 $\ln2<1$，确实比远离时小。两球越接近，同号电荷越被排向外侧，
同样总电荷产生的公共电势越高，所以 $C=Q/V$ 越小，趋势与级数结果一致。

“两球导通”是本题的关键条件。若两球仅几何接触却由绝缘层隔开，电荷不能在两个球之间重新分配，
就不能把表面视为同一个等势导体；若两球带给定的不同总电荷，也应先求电容系数矩阵，而不是只用
一个标量电容。图示中的连成一体，既表示几何接触，也表示电荷能够自由迁移并达到共同电势。
\end{xesolution}

\Needspace{0.18\textheight}
\begin{origexamplebox}{例题 32.3【两个圆柱】}
两根无穷长平行圆柱导体的半径均为 $R$，圆心距为 $L>2R$。两柱带等量异号线电荷，求单位
长度的双极电容。
\end{origexamplebox}

\begin{xesolution}
取双极坐标 $\tau,\sigma$，使两圆柱表面分别为 $\tau=\pm\tau_0$。几何关系为
\begin{equation}
  \boxed{\cosh\tau_0=\frac{L}{2R}}.
\end{equation}
二维 Laplace 方程中，满足两柱电势分别为 $+U/2$、$-U/2$ 的解可取
\[
  \varphi(\tau)=\frac{U}{2\tau_0}\tau.
\]
由柱面上的法向电场积分，单根圆柱每单位长度所带电荷为
\[
  \lambda=\frac{\pi\varepsilon_0U}{\tau_0}.
\]
因此单位长度双极电容
\begin{equation}
  \boxed{C'=\frac\lambda U
  =\frac{\pi\varepsilon_0}{\operatorname{arcosh}(L/2R)}}.
\end{equation}
利用 $\operatorname{arcosh}x=\ln(x+\sqrt{x^2-1})$，也可写成
\begin{equation}
  \boxed{C'=\frac{\pi\varepsilon_0}
  {\ln\!\left(\dfrac{L+\sqrt{L^2-4R^2}}{2R}\right)}}.
\end{equation}
$L\to2R^+$ 时两柱接近接触，电容发散；这是理想无限长、零间隙模型的结果。

\twoCylinderDiagram
\end{xesolution}

\Needspace{0.4\textheight}
\begin{origexamplebox}{例题 32.4【改造电容器】}
平行板电容器的极板面积为 $S$、间距为 $d$。
\begin{enumerate}[label=(\arabic*),leftmargin=2.2em]
  \item 沿面积方向，在左半边完全填入相对介电常数为 $\varepsilon_r$ 的介质；
  \item 沿厚度方向，在上半层填入相对介电常数为 $\varepsilon_r$ 的介质。
\end{enumerate}
分别求改造后的电容，并判断增减。
\end{origexamplebox}

\begin{xesolution}
\textbf{(1) 按面积分成左右两区。}两区具有相同板间距与电压，相当于并联：
\[
  C_1=\frac{\varepsilon_r\varepsilon_0(S/2)}d,
  \qquad C_2=\frac{\varepsilon_0(S/2)}d.
\]
故
\begin{equation}
  \boxed{C_{\rm left}=C_1+C_2
  =\frac{\varepsilon_0S}{2d}(\varepsilon_r+1)}.
\end{equation}

\textbf{(2) 按厚度分成上下两层。}两层面积相同，自由面电荷密度相同，电位移的法向分量连续，
相当于串联：
\[
  C_1=\frac{\varepsilon_r\varepsilon_0S}{d/2},
  \qquad C_2=\frac{\varepsilon_0S}{d/2}.
\]
于是
\begin{equation}
  \boxed{C_{\rm upper}=\left(\frac1{C_1}+\frac1{C_2}\right)^{-1}
  =\frac{2\varepsilon_r}{\varepsilon_r+1}\frac{\varepsilon_0S}{d}}.
\end{equation}
因 $\varepsilon_r>1$，两种改造都使电容变大；但一般有
\[
  C_{\rm left}-C_{\rm upper}
  =\frac{\varepsilon_0S}{d}\frac{(\varepsilon_r-1)^2}{2(\varepsilon_r+1)}\ge0,
\]
所以相同介质量按面积并排放置，增容更明显。

\splitDielectricDiagram
\end{xesolution}

\Needspace{0.42\textheight}
\begin{origexamplebox}{例题 32.5【奇妙的吸引力】}
小明又需要你的帮助了。他把一块相对介电常数为 $\varepsilon_r$ 的电介质插入平行板电容器中。
极板面积为 $S=ab$、间距为 $d$，介质厚度等于板间距，并沿边长 $a$ 的方向插入深度 $x$。
电介质会受到继续向电容器内部运动的吸引力；这个力从微观上怎样产生并不容易直接看清，请利用虚功法
把它算出来。
\begin{enumerate}[label=(\arabic*),leftmargin=2.2em]
  \item 电容器电荷 $Q$ 保持不变，用介质内电场 $E$ 表示吸入力；
  \item 电容器接在恒压电源上，用介质内电场 $E$ 表示吸入力。
\end{enumerate}
\end{origexamplebox}

\begin{xesolution}
忽略边缘场，已填充区域与真空区域相当于并联：
\begin{equation}
  C(x)=\frac{\varepsilon_0b}{d}\,[a+(\varepsilon_r-1)x],
  \qquad
  \frac{\dd C}{\dd x}=\frac{\varepsilon_0b}{d}(\varepsilon_r-1).
\end{equation}

\textbf{(1) 固定电荷。}系统静电能 $W=Q^2/[2C(x)]$，介质所受力沿 $x$ 增大方向：
\[
  F=-\left(\frac{\partial W}{\partial x}\right)_Q
  =\frac{Q^2}{2C^2}\frac{\dd C}{\dd x}
  =\frac{U^2}{2}\frac{\dd C}{\dd x}.
\]
介质内场 $E=U/d$，故
\begin{equation}
  \boxed{F_Q=\frac12\varepsilon_0(\varepsilon_r-1)bdE^2}.
\end{equation}

\textbf{(2) 固定电压。}电容器能量 $CU^2/2$ 随 $C$ 增大而增大，但电源还会输送电荷并做功
$U\dd Q=U^2\dd C$。把电源计入后，机械势能的增量是
$\dd(CU^2/2)-U\dd Q=-U^2\dd C/2$，所以
\[
  F=\frac{U^2}{2}\frac{\dd C}{\dd x}.
\]
仍用 $E=U/d$ 表示，得到
\begin{equation}
  \boxed{F_U=\frac12\varepsilon_0(\varepsilon_r-1)bdE^2}.
\end{equation}
两式形式相同，是因为题目要求用各自状态下的 $E$ 表示；固定 $Q$ 时 $E$ 随 $x$ 改变，固定
$U$ 时 $E$ 不变。力都把介质拉入电容器，因为介质进入会增大电容。

\dielectricPullDiagram
\end{xesolution}

\Needspace{0.4\textheight}
\begin{origexamplebox}{例题 32.6【匀强电场中的电介质球】}
半径为 $R$、介电常数为 $\varepsilon$ 的均匀线性电介质球置于真空匀强电场 $\vect E_0$ 中，求
空间中的电势与电场。
\end{origexamplebox}

\begin{xesolution}
取 $z$ 轴沿 $\vect E_0$。球内、外均无自由电荷，电势满足 Laplace 方程。只保留满足轴对称和
无穷远条件的一阶项：
\[
  \varphi_{\rm in}=-Ar\cos\theta,
  \qquad
  \varphi_{\rm out}=-E_0r\cos\theta+\frac{B\cos\theta}{r^2}.
\]
在 $r=R$ 处电势连续，切向电场自然连续；又因界面无自由面电荷，法向电位移连续：
\[
  \varphi_{\rm in}=\varphi_{\rm out},
  \qquad
  \varepsilon E_{r,\rm in}=\varepsilon_0E_{r,\rm out}.
\]
解得
\[
  A=\frac{3\varepsilon_0}{\varepsilon+2\varepsilon_0}E_0,
  \qquad
  B=\frac{\varepsilon-\varepsilon_0}{\varepsilon+2\varepsilon_0}E_0R^3.
\]
因此
\begin{equation}
  \boxed{\varphi_{\rm in}=-\frac{3\varepsilon_0}{\varepsilon+2\varepsilon_0}
  E_0r\cos\theta},
  \qquad
  \boxed{\vect E_{\rm in}=\frac{3\varepsilon_0}{\varepsilon+2\varepsilon_0}\vect E_0}.
\end{equation}
球内场仍为匀强场，但弱于外加场。球外
\begin{equation}
  \boxed{\varphi_{\rm out}=-E_0r\cos\theta
  +\frac{\varepsilon-\varepsilon_0}{\varepsilon+2\varepsilon_0}
  E_0R^3\frac{\cos\theta}{r^2}},
  \qquad \vect E_{\rm out}=-\nabla\varphi_{\rm out}.
\end{equation}
第二项是感应偶极子的电势，对应
\begin{equation}
  \boxed{\vect p=4\pi\varepsilon_0R^3
  \frac{\varepsilon-\varepsilon_0}{\varepsilon+2\varepsilon_0}\vect E_0}.
\end{equation}
当 $\varepsilon\to\infty$ 时，球内场趋于零，外部解退化为上一讲的导体球结果。

\dielectricSphereDiagram
\end{xesolution}

\Needspace{0.34\textheight}
\begin{origexamplebox}{例题 32.7【能量怎么多了】}
线性电介质中的能量密度为 $u=\varepsilon E^2/2$，真空中则为
$u_0=\varepsilon_0E^2/2$。在电场强度相同的条件下，介质中的能量为什么更多？
\end{origexamplebox}

\begin{xesolution}
介质中的总能量不只包括真空电场本身的能量，还包括建立极化所需的微观形变能。线性介质中
\[
  \vect D=\varepsilon_0\vect E+\vect P.
\]
从零场准静态建立到 $\vect E$，单位体积外界输入
\[
  u=\int_{\vect 0}^{\vect D}\vect E\cdot\dd\vect D
  =\frac12\vect E\cdot\vect D
  =\frac12\varepsilon_0E^2+\frac12\vect E\cdot\vect P.
\]
第一项是真空场能；第二项是把束缚正、负电荷拉开，建立分子偶极矩与介质极化所需要储存的能量。
因此在比较相同 $E$ 时，
\begin{equation}
  \boxed{u-u_0=\frac12\vect E\cdot\vect P
  =\frac12(\varepsilon-\varepsilon_0)E^2>0}.
\end{equation}
若比较的是固定自由电荷而不是固定 $E$，插入介质会使 $E$ 减小，总能量反而可能下降；此时减少
的电场能转化为介质被吸入时的机械功。比较能量前必须先明确保持不变的是 $E$、$Q$ 还是 $U$。
\end{xesolution}

\Needspace{0.4\textheight}
\begin{origexamplebox}{例题 32.8【充电过程】}
用电动势为 $\mathcal E$ 的理想电源给电容 $C$ 充电。
\begin{enumerate}[label=(\arabic*),leftmargin=2.2em]
  \item 电路中串联电阻 $R$，充电完成前电源做功和电阻发热各为多少；
  \item 令 $R$ 趋于零，两个结果怎样变化；
  \item 若理想电路中完全没有电阻，应如何理解？
\end{enumerate}
\end{origexamplebox}

\begin{xesolution}
\textbf{(1)} 充电完成后 $Q=C\mathcal E$。电源搬运总电荷 $Q$，故
\begin{equation}
  \boxed{W_{\rm source}=\mathcal EQ=C\mathcal E^2}.
\end{equation}
电容器最终储能
\[
  W_C=\frac12C\mathcal E^2.
\]
其余能量变成焦耳热：
\begin{equation}
  \boxed{Q_R=W_{\rm source}-W_C=\frac12C\mathcal E^2}.
\end{equation}
也可由 $i(t)=\mathcal E e^{-t/(RC)}/R$ 积分验证：
$\int_0^\infty i^2R\dd t=C\mathcal E^2/2$。

\textbf{(2)} 当 $R\to0^+$，充电时间常数 $RC\to0$、初始电流峰值趋于无穷，但热量积分仍为
$C\mathcal E^2/2$，与 $R$ 的数值无关。这是先对有限 $R$ 完成积分、再取极限的结果。

\textbf{(3)} 完全没有电阻、电感和辐射的理想回路不能由一阶 $RC$ 方程描述；它要求电容电压在
零时间内突变，产生无穷大电流，数学模型自相矛盾。真实电路总有寄生电感、导线电阻和电磁辐射。
若保留电感而不含耗散，能量会在电容电场与电感磁场之间振荡，电容电压并不单调趋于
$\mathcal E$；微小电阻或辐射最终耗散那一半能量。因而“零电阻”只能作为 $R\to0^+$ 的极限，
不能把所有寄生参数同时删掉后仍期待同样的稳态充电过程。

\chargingDiagram
\end{xesolution}

\Needspace{0.38\textheight}
\begin{origexamplebox}{例题 32.9【电容充电】}
三个初始均不带电的电容 $C_1,C_2,C_3$ 连接成星形，外端电势分别固定为
$\varphi_1,\varphi_2,\varphi_3$，中心节点悬空。求中心电势及三个电容器所带电荷。
\end{origexamplebox}

\begin{xesolution}
设中心节点电势为 $\varphi_0$。中心节点及与它相连的三块极板构成一个孤立导体，初始净电荷为
零，此后也没有导电通路让净电荷流入，因此
\[
  C_1(\varphi_0-\varphi_1)+C_2(\varphi_0-\varphi_2)
  +C_3(\varphi_0-\varphi_3)=0.
\]
解得中心电势
\begin{equation}
  \boxed{\varphi_0=
  \frac{C_1\varphi_1+C_2\varphi_2+C_3\varphi_3}{C_1+C_2+C_3}}.
\end{equation}
若把 $Q_i$ 定义为第 $i$ 个外端极板的电荷，则
\begin{equation}
  \boxed{Q_i=C_i(\varphi_i-\varphi_0)},\qquad i=1,2,3.
\end{equation}
展开为
\begin{align*}
  Q_1&=\frac{C_1[C_2(\varphi_1-\varphi_2)+C_3(\varphi_1-\varphi_3)]}{C_1+C_2+C_3},\\
  Q_2&=\frac{C_2[C_1(\varphi_2-\varphi_1)+C_3(\varphi_2-\varphi_3)]}{C_1+C_2+C_3},\\
  Q_3&=\frac{C_3[C_1(\varphi_3-\varphi_1)+C_2(\varphi_3-\varphi_2)]}{C_1+C_2+C_3}.
\end{align*}
中心侧极板分别带 $-Q_i$，且 $Q_1+Q_2+Q_3=0$，正好体现中心孤立导体的电荷守恒。

\starCapacitorDiagram
\end{xesolution}

\Needspace{0.34\textheight}
\begin{origexamplebox}{例题 32.10【电容串并联】}
证明两个电容器串联或并联都可等效为一个电容器，求等效公式，并推广到多个电容器。
\end{origexamplebox}

\begin{xesolution}
\textbf{并联。}各电容电压相同，均为 $U$，电源输送的总电荷为
\[
  Q=\sum_iQ_i=\sum_iC_iU.
\]
令 $Q=C_{\rm eq}U$，得到
\begin{equation}
  \boxed{C_{\rm eq}=\sum_iC_i}.
\end{equation}

\textbf{串联。}各中间节点初始中性且保持孤立，所以每个电容所带电荷量的绝对值均为 $Q$。总
电压为
\[
  U=\sum_iU_i=Q\sum_i\frac1{C_i}.
\]
由 $Q=C_{\rm eq}U$ 得
\begin{equation}
  \boxed{\frac1{C_{\rm eq}}=\sum_i\frac1{C_i}}.
\end{equation}
两个电容时，$C_{\rm series}=C_1C_2/(C_1+C_2)$，
$C_{\parallel}=C_1+C_2$。等效不仅保持端口处的 $Q$-$U$ 关系，也保持端口看到的总储能：
\[
  \sum_i\frac12C_iU_i^2=\frac12C_{\rm eq}U^2.
\]
若中间节点预先带有非零净电荷，或存在额外支路，简单“串联同电荷”条件不再成立，需要按节点
电荷守恒重新求解。
\end{xesolution}

\Needspace{0.44\textheight}
\begin{origexamplebox}{例题 32.11【超大电容】}
物理学家 Kelvin 最先把地球和大气层构造为一个电容器模型：地球表面是其中一个极板，带负电荷；
高空的大气等效为另一个极板，带正电荷。把地球表面和高空 $h=\SI{5}{km}$ 处的大气层近似为
同心球形电容器。取地球半径
$R_\oplus=\SI{6.4e6}{m}$，地球带电量绝对值 $Q=\SI{5e5}{C}$，空气电阻率
$\rho=\SI{3e13}{\ohm.m}$。
\begin{enumerate}[label=(\arabic*),leftmargin=2.2em]
  \item 求球形电容器的电容；
  \item 求地表电场能量密度和电容器总能量；
  \item 求大气层的径向漏电电阻；
  \item 求 $RC$ 放电时间常数；
  \item 若每次雷电平均向地面补充 $\SI{25}{C}$ 电荷，维持该电荷量每天约需多少次雷电？
\end{enumerate}
\end{origexamplebox}

\begin{xesolution}
令内、外半径分别为 $R=R_\oplus$、$R+h$。

\textbf{(1) 电容。}
\begin{equation}
  C=4\pi\varepsilon_0\frac{R(R+h)}h
  \approx\boxed{\SI{0.912}{F}}.
\end{equation}
因 $h\ll R$，也可用薄层近似 $C\approx4\pi\varepsilon_0R^2/h$。

\textbf{(2) 地表场与能量。}地表外侧
\[
  E(R)=\frac{Q}{4\pi\varepsilon_0R^2}approx\SI{1.10e2}{V.m^{-1}}.
\]
故
\begin{equation}
  \boxed{u(R)=\frac12\varepsilon_0E^2\approx\SI{5.33e-8}{J.m^{-3}}}.
\end{equation}
总能量
\begin{equation}
  \boxed{W=\frac{Q^2}{2C}\approx\SI{1.37e11}{J}}.
\end{equation}

\textbf{(3) 径向电阻。}半径 $r$ 处厚度 $\dd r$ 的球壳电阻
$\dd R_{\rm el}=\rho\dd r/(4\pi r^2)$，所以
\begin{equation}
  R_{\rm el}=\frac\rho{4\pi}\left(\frac1R-\frac1{R+h}\right)
  \approx\boxed{\SI{2.91e2}{\ohm}}.
\end{equation}

\textbf{(4) 时间常数。}相乘时几何因子恰好消去：
\begin{equation}
  \boxed{\tau=R_{\rm el}C=\rho\varepsilon_0
  \approx\SI{2.66e2}{s}\approx\SI{4.43}{min}}.
\end{equation}

\textbf{(5) 雷电补充率。}漏电电流的平均量级为
\[
  I=\frac Q\tau\approx\SI{1.88e3}{A}.
\]
每天漏失电荷 $I\times\SI{86400}{s}$，所需雷电次数
\begin{equation}
  \boxed{N=\frac{Q}{\tau}\frac{\SI{86400}{s}}{\SI{25}{C}}
  \approx6.5\times10^6\ \text{次/天}}.
\end{equation}
这是题设参数下的数量级模型；真实全球电路的电荷分布、天气电导率和雷暴补充机制更复杂。

\atmosphereCapacitorDiagram
\end{xesolution}
